\par\Needspace{6\baselineskip}\medskip
\hypertarget{native-151-article109}{}\pdfbookmark[2]{109.}{article-109}
\begin{center}109.\end{center}
If we use the symbols $\alpha$, $C$, and $c$ for the hyperbola in the same sense as in Article 99, equations VIII and IX of Article 21 give
\[
\begin{aligned}
r'\cos v'-r\cos v&=-\tfrac12\left(c-\frac1c\right)\left(C-\frac1C\right)\alpha\\[.7em]
r'\sin v'-r\sin v&=\tfrac12\left(c-\frac1c\right)\left(C+\frac1C\right)\alpha\BareRoot{(ee-1)},
\end{aligned}
\]
\noindent and hence
\[
\rho=\tfrac12\alpha\left(c-\frac1c\right)\BareRoot{\left(ee\left(C+\frac1C\right)^2-4\right)}.
\]
\noindent\SourceLeaf{152}Let $\gamma$ be the quantity determined by the equation $\gamma+\dfrac1\gamma=e\left(C+\dfrac1C\right)$. Since it is clear that \emph{two} mutually reciprocal values satisfy it, we take the one greater than 1. This gives
\[
\rho=\tfrac12\alpha\left(c-\frac1c\right)\left(\gamma-\frac1\gamma\right)\mathpunct{\cdot}
\]
\noindent Furthermore,
\[
\begin{aligned}
r+r'&=\tfrac12\alpha\left(e\left(c+\frac1c\right)\left(C+\frac1C\right)-4\right)\\[.7em]
&=\tfrac12\alpha\left(\left(c+\frac1c\right)\left(\gamma+\frac1\gamma\right)-4\right),
\end{aligned}
\]
\noindent and hence
\[
\begin{aligned}
r+r'+\rho&=\alpha\left(\BareRoot{c\gamma}-\BareRoot{\frac1{c\gamma}}\right)^2\\[.7em]
r+r'-\rho&=\alpha\left(\BareRoot{\frac\gamma c}-\BareRoot{\frac c\gamma}\right)^2\mathpunct{\cdot}
\end{aligned}
\]
\noindent If we therefore set
\[
\BareRoot{\frac{r+r'+\rho}{4\alpha}}=m,\qquad\BareRoot{\frac{r+r'-\rho}{4\alpha}}=n,
\]
\noindent we must have $\BareRoot{c\gamma}-\BareRoot{\dfrac1{c\gamma}}=2m$. To decide whether $\BareRoot{\dfrac\gamma c}-\BareRoot{\dfrac c\gamma}$ is $=+2n$ or $=-2n$, we must determine whether $\gamma$ is greater or less than $c$. It follows readily from equation 8 of Article 99 that the former case occurs whenever $2f$ is below $180^0$ and the latter whenever $2f$ is above $180^0$. Finally, from the same article we have
\[
\begin{aligned}
\frac{kt}{\alpha^{\frac32}}&=\tfrac12\left(\gamma+\frac1\gamma\right)\left(c-\frac1c\right)-2\log c\\[.8em]
&=\tfrac12\left(c\gamma-\frac1{c\gamma}\right)-\tfrac12\left(\frac\gamma c-\frac c\gamma\right)-\log c\gamma+\log\frac\gamma c\\[.8em]
&=2m\BareRoot{(1+mm)}\mp2n\BareRoot{(1+nn)}\\[.5em]
&\qquad-2\log\bigl(\BareRoot{(1+mm)}+m\bigr)\pm2\log\bigl(\BareRoot{(1+nn)}+n\bigr)
\end{aligned}
\]
\noindent where the lower signs always apply to the case $2f>180^0$. Now
\[
\log\bigl(\BareRoot{(1+mm)}+m\bigr)
\]
\noindent readily expands into the following series
\[
m-\tfrac13\cdot\tfrac12m^3+\tfrac15\cdot\frac{1\mathpunct{.}3}{2\mathpunct{.}4}m^5-\tfrac17\cdot\frac{1\mathpunct{.}3\mathpunct{.}5}{2\mathpunct{.}4\mathpunct{.}6}m^7+\text{etc.}
\]
\noindent This follows directly from $\mathrm d\log\bigl(\BareRoot{(1+mm)}+m\bigr)=\dfrac{\mathrm dm}{\BareRoot{(1+mm)}}\mathpunct{\cdot}$ We thus obtain
\[
\begin{aligned}
&2m\BareRoot{(1+mm)}-2\log\bigl(\BareRoot{(1+mm)}+m\bigr)\\[.6em]
&\qquad=4\left(\tfrac13m^3-\tfrac15\cdot\tfrac12m^5+\tfrac17\cdot\frac{1\mathpunct{.}3}{2\mathpunct{.}4}m^7-\text{etc.}\right)
\end{aligned}
\]
\noindent\SourceLeaf{153}and likewise an entirely similar formula if $m$ is interchanged with $n$. Finally, if we set
\[
\begin{aligned}
T={}&\tfrac16(r+r'-\rho)^{\frac32}-\tfrac1{80}\cdot\frac1\alpha(r+r'-\rho)^{\frac52}\\[.6em]
&+\tfrac3{1792}\cdot\frac1{\alpha\alpha}(r+r'-\rho)^{\frac72}-\tfrac5{18432}\cdot\frac1{\alpha^3}(r+r'-\rho)^{\frac92}+\text{etc.}
\end{aligned}
\]
\[
\begin{aligned}
U={}&\tfrac16(r+r'+\rho)^{\frac32}-\tfrac1{80}\cdot\frac1\alpha(r+r'+\rho)^{\frac52}\\[.6em]
&+\tfrac3{1792}\cdot\frac1{\alpha\alpha}(r+r'+\rho)^{\frac72}-\tfrac5{18432}\cdot\frac1{\alpha^3}(r+r'+\rho)^{\frac92}+\text{etc.,}
\end{aligned}
\]
\noindent it follows that $kt=U\mp T$. These expressions coincide exactly with those given in Article 107 if $a$ there is replaced by $-\alpha$.
\par
For both the ellipse and the hyperbola, these series are especially convenient in practice when $a$ or $\alpha$ takes a very large value, that is, when the conic section closely resembles a parabola. In such a case they could also be used to solve the problem treated above (Articles 85--105). Since in our judgment they do not offer the brevity of the solution given above even then, we do not dwell on a fuller exposition of this method.
\par\bigskip\begin{center}\rule{.23\linewidth}{.35pt}\end{center}
